\documentclass[12pt,a4paper,sumario=tradicional]{abntex2}

% --- Pacotes ---
\usepackage[utf8]{inputenc}
\usepackage[T1]{fontenc}
% babel é carregado automaticamente pelo abntex2
\usepackage{lastpage}
\usepackage{graphicx}
\usepackage{xcolor}
\usepackage{hyperref}
\hypersetup{colorlinks=true, linkcolor=blue, citecolor=blue, urlcolor=blue}
\usepackage{booktabs}
\usepackage{longtable}
\usepackage{multirow}
\usepackage{listings}
\usepackage{amsmath, amssymb}
\usepackage{caption}
\usepackage{times}
\usepackage[alf]{abntex2cite}

% --- Configurações ---
\lstset{
    basicstyle=\footnotesize\ttfamily,
    breaklines=true,
    frame=single,
    numbers=left,
    numberstyle=\tiny,
    backgroundcolor=\color{gray!5}
}

% --- Metadados ABNT ---
\titulo{EVOLUÇÃO DO ECOSSISTEMA OPENCODE}
\autor{Marcelo Claro Laranjeira}
\orientador{[ORIENTADOR]}
\coorientador{[COORIENTADOR]}
\instituicao{Universidade Federal do Ceará}
\local{Fortaleza -- Ceará}
\data{2026}
\tipotrabalho{dissertacao}
\preambulo{Dissertação apresentada ao Programa de Pós-Graduação em Tecnologia Educacional da Universidade Federal do Ceará como requisito parcial para obtenção do título de Mestre em Tecnologia Educacional.}

\begin{document}

% ======================================================================
% CAPA
% ======================================================================
\imprimircapa

% ======================================================================
% FOLHA DE ROSTO
% ======================================================================
\imprimirfolhaderosto

% ======================================================================
% FICHA CATALOGRÁFICA
% ======================================================================
\newpage
\thispagestyle{empty}
\vspace*{10cm}
\begin{center}
\fbox{\parbox{0.85\textwidth}{
\centering
\textbf{Ficha Catalográfica}\\[0.5cm]
LARANJEIRA, Marcelo Claro\\
Evolução do Ecossistema OpenCode: do raciocínio frágil à elegância autônoma --\\
Uma trajetória de 13 ciclos evolutivos com 212 tipos de raciocínio e 125 agentes.\\
Fortaleza, 2026.\\
\pageref{LastPage} p.\\[0.3cm]
Dissertação (Mestrado Profissional em Tecnologia Educacional) --\\
Universidade Federal do Ceará, Programa de Pós-Graduação em Tecnologia Educacional.\\[0.3cm]
1. Inteligência Artificial. 2. Sistemas Multiagente. 3. Ecossistema Evolutivo.\\
4. Verificação Formal. 5. OpenCode. I. Título.
}}
\end{center}

% ======================================================================
% AGRADECIMENTOS
% ======================================================================
\newpage
\chapter*{Agradecimentos}
\addcontentsline{toc}{chapter}{Agradecimentos}

Ao Programa de Pós-Graduação em Tecnologia Educacional (PPGTE/UFC) pelo ambiente de pesquisa e inovação.

Aos orientadores e colegas do GeoMaker+IA/CT que contribuíram com discussões e revisões ao longo dos 13 ciclos evolutivos.

À comunidade open-source dos projetos MiroFish, BettaFish, DeerFlow, Agent Skills, Open Design e demais repositórios que inspiraram padrões arquiteturais fundamentais.

Aos 125 agentes do ecossistema, cujo trabalho incansável demonstra que a inteligência artificial pode produzir ciência verificável, ética e de qualidade Qualis A1.

% ======================================================================
% RESUMO
% ======================================================================
\newpage
\begin{center}
\textbf{RESUMO}
\end{center}

Este trabalho documenta a evolução completa do Ecossistema OpenCode, uma plataforma multiagente autônoma que evoluiu de um sistema de raciocínio frágil para uma arquitetura de elegância computacional autônoma ao longo de 13 ciclos evolutivos. A trajetória abrange desde a primeira geração (Evo-1 a Evo-6: análise manual de correlações econômicas, pipeline de artigos acadêmicos, sistema TSAC de citações) até a terceira geração (Evo-10 a Evo-13: DataOrchestrator multi-domínio, Reasoning Orchestrator com 68 tipos, Auditoria Caixa Branca). O ecossistema integra 125 agentes especializados, 40 MCPs (Model Context Protocols), 106 skills em 13 categorias, 15 plugins e 14 comandos slash, todos orquestrados por um pipeline de verificação formal Cora-Debate com 6 verificadores simbólicos (V1-V6). São documentados 18 padrões arquiteturais extraídos dos projetos MiroFish e BettaFish (P1-P18), o sistema PhD Auditor com 10 critérios de qualidade Qualis A1, a integração com 52 editais de fomento brasileiros (editais-br v7.1) e a validação TDD de 5 Especificações Técnicas (SPECs) com 25/25 testes aprovados. A plataforma alcançou Health Score 100/100, 212 tipos de raciocínio em 27 categorias e 100\% de aprovação nos testes funcionais do Cora-Debate.

\vspace{0.5cm}
\textbf{Palavras-chave}: Inteligência Artificial. Sistemas Multiagente. Ecossistema Evolutivo. Verificação Formal. OpenCode.

% ======================================================================
% ABSTRACT
% ======================================================================
\newpage
\begin{center}
\textbf{ABSTRACT}
\end{center}

This work documents the complete evolution of the OpenCode Ecosystem, an autonomous multi-agent platform that evolved from a fragile reasoning system into an autonomous computational elegance architecture over 13 evolutionary cycles. The trajectory spans from the first generation (Evo-1 to Evo-6: manual economic correlation analysis, academic article pipeline, TSAC citation system) to the third generation (Evo-10 to Evo-13: multi-domain DataOrchestrator, Reasoning Orchestrator with 68 types, White Box Audit). The ecosystem integrates 125 specialized agents, 40 MCPs (Model Context Protocols), 106 skills in 13 categories, 15 plugins and 14 slash commands, all orchestrated by a Cora-Debate formal verification pipeline with 6 symbolic verifiers (V1-V6). We document 18 architectural patterns extracted from MiroFish and BettaFish projects (P1-P18), the PhD Auditor system with 10 Qualis A1 quality criteria, integration with 52 Brazilian funding calls (editais-br v7.1), and TDD validation of 5 Technical Specifications (SPECs) with 25/25 passing tests. The platform achieved Health Score 100/100, 212 reasoning types in 27 categories and 100\% approval in Cora-Debate functional tests.

\vspace{0.5cm}
\textbf{Keywords}: Artificial Intelligence. Multiagent Systems. Evolutionary Ecosystem. Formal Verification. OpenCode.

% ======================================================================
% LISTA DE FIGURAS E TABELAS
% ======================================================================
\newpage
\listoffigures
\vspace{1cm}
\listoftables

% ======================================================================
% SUMÁRIO
% ======================================================================
\newpage
\tableofcontents
\newpage

% ======================================================================
% CAPÍTULO 1: INTRODUÇÃO
% ======================================================================
\chapter{Introdução}

\section{Contexto}

O OpenCode nasceu como um assistente de codificação para terminal, inspirado pelo Claude Code da Anthropic. No entanto, sua arquitetura modular -- baseada em Model Context Protocols (MCPs), skills carregáveis sob demanda e plugins TypeScript -- permitiu uma expansão que transcendeu o propósito original. Em 13 ciclos evolutivos documentados, o sistema evoluiu de um analisador de correlações econômicas para um ecossistema multiagente completo com verificação formal, produção acadêmica Qualis A1, computação quântica simulada e engenharia reversa autônoma.

\section{Problema de Pesquisa}

Como projetar um ecossistema de agentes de inteligência artificial que:
\begin{enumerate}
    \item Evolua autonomamente sem intervenção humana;
    \item Produza raciocínio científico verificável com garantias formais;
    \item Integre múltiplos paradigmas (lógico, dialético, teoria dos jogos, estatístico);
    \item Mantenha rastreabilidade completa de decisões e custos;
    \item Alcance padrão de qualidade acadêmica Qualis A1.
\end{enumerate}

\section{Objetivos}

\subsection{Objetivo Geral}

Documentar e analisar a evolução do Ecossistema OpenCode como plataforma multiagente autônoma para produção de conhecimento científico verificável.

\subsection{Objetivos Específicos}

\begin{enumerate}
    \item Sistematizar os 13 ciclos evolutivos em três gerações;
    \item Analisar os 18 padrões arquiteturais integrados (P1-P18);
    \item Validar a arquitetura Cora-Debate com 6 verificadores simbólicos;
    \item Demonstrar a aplicação do pipeline TDD acadêmico com 5 SPECs;
    \item Avaliar o sistema PhD Auditor com 10 critérios Qualis A1.
\end{enumerate}

\section{Estrutura do Documento}

O Capítulo 2 apresenta os fundamentos teóricos. O Capítulo 3 documenta a primeira geração evolutiva (Evo-1 a Evo-6). O Capítulo 4 descreve a segunda geração com sincronização autônoma (Evo-7 a Evo-9). O Capítulo 5 aborda a terceira geração com orquestração multiagente (Evo-10 a Evo-13). O Capítulo 6 detalha a arquitetura Cora-Debate. O Capítulo 7 apresenta o pipeline de Engenharia Reversa. O Capítulo 8 documenta a integração MiroFish/BettaFish. O Capítulo 9 descreve o sistema TDD Acadêmico com SPECs. O Capítulo 10 apresenta o PhD Auditor. O Capítulo 11 conclui com métricas agregadas e perspectivas futuras.

% ======================================================================
% CAPÍTULO 2: FUNDAMENTOS
% ======================================================================
\chapter{Fundamentos Teóricos}

\section{Sistemas Multiagente}

A teoria de sistemas multiagente, conforme Wooldridge (2009), estabelece que agentes inteligentes podem cooperar, competir e negociar em ambientes compartilhados. O Ecossistema OpenCode estende esse paradigma ao introduzir evolução autônoma: os agentes não apenas executam tarefas, mas criam novos skills, plugins e padrões de raciocínio.

\section{Verificação Formal}

A verificação formal de raciocínio científico combina:
\begin{itemize}
    \item \textbf{Análise dimensional:} verificação de consistência de unidades físicas;
    \item \textbf{Verificação algébrica:} simplificação simbólica via SymPy;
    \item \textbf{Geração de contraexemplos:} busca randomizada de casos limite;
    \item \textbf{Validação estatística:} testes de hipótese com correção de Bonferroni;
    \item \textbf{Verificação numérica:} precisão IEEE 754 float64;
    \item \textbf{Validação de EDOs/PDEs:} solução simbólica via SymPy dsolve.
\end{itemize}

\section{Ecossistemas Evolutivos}

O conceito de ``ecossistema evolutivo'' aplicado a plataformas de IA refere-se à capacidade de auto-modificação, auto-correção e auto-expansão sem intervenção humana. O ciclo PLAN$\rightarrow$ACT$\rightarrow$REFLECT$\rightarrow$EXTRACT$\rightarrow$EVOLVE, implementado pelo plugin \texttt{manus-evolve.ts}, permitiu 13 gerações de evolução com score crescente (85$\to$100).

\section{Process Confidence Index (PCI)}

O PCI é uma métrica que quantifica a confiança em raciocínios gerados por IA. A fórmula de calibração PCI (SPEC-PCI-001) ajusta o valor bruto por:
\begin{equation}
    \text{PCI}_{\text{cal}} = \min\left(10, \max\left(0, \text{PCI}_{\text{bruto}} \times 0.062 \times f_{\text{domínio}} + b_{\text{código}} + b_{\text{narração}}\right)\right)
\end{equation}

onde $f_{\text{domínio}} \in \{1.00, 0.85, 0.70\}$ para domínios geral, geometria e numérico, respectivamente.

% ======================================================================
% CAPÍTULO 3: PRIMEIRA GERAÇÃO (EVO-1 A EVO-6)
% ======================================================================
\chapter{Primeira Geração Evolutiva: Análise Manual}

\section{Evo-1: Armadilha da Renda Média (Score: 85/100)}

O primeiro ciclo evolutivo focou na análise da ``Armadilha da Renda Média'' (Middle Income Trap) utilizando 27 indicadores do Banco Mundial para 10 países.

\textbf{Principais descobertas:}
\begin{itemize}
    \item Correlação educação vs PIB per capita: $r = -0.03$ (não significativa);
    \item Correlação P\&D privado vs PIB per capita: $r = +0.73$ (forte);
    \item 5 ferramentas integradas em pipeline unificado.
\end{itemize}

\textbf{Métrica:} 27 indicadores, 10 países, correlação Pearson bivariada.

\section{Evo-2: Pipeline de Artigo Acadêmico (Score: 90/100)}

O segundo ciclo expandiu a capacidade de produção acadêmica:
\begin{itemize}
    \item 11 arquivos gerados, 11.476 palavras, 11 tabelas;
    \item Descoberta principal: Serviços de Alta Tecnologia como preditor $r = 0.95$;
    \item Pipeline artigo acadêmico com escrita anti-IA.
\end{itemize}

\section{Evo-3: Sistema TSAC de Citações + Sci-Hub (Score: 92/100)}

\begin{itemize}
    \item 46 notas TSAC (Textual Source Analysis Credibility);
    \item Cada nota com 5 componentes: referência ABNT + DOI + trecho original + tradução + justificativa;
    \item Pipeline Sci-Hub MCP + arXiv API + Semantic Scholar + fallback;
    \item Integração com CrossRef para validação de DOIs.
\end{itemize}

\section{Evo-4/5: Cross-Validation Engine (Score: 92/100)}

\begin{itemize}
    \item Correlação Pearson em 27 indicadores econômicos;
    \item 3 níveis de validação cruzada;
    \item Auto-scoring com critérios Qualis A1.
\end{itemize}

\section{Evo-6: Iterative Correction Loop (Score: 95/100)}

O ciclo mais intensivo da primeira geração:

\begin{longtable}{@{}lcc@{}}
\caption{Evo-6: Métricas de Correção Iterativa}\\
\toprule
\textbf{Métrica} & \textbf{Antes} & \textbf{Depois} \\
\midrule
Board Score & 86.5 & 92.7 \\
Auto Score Qualis & 74 & 95 \\
DOIs verificados & 12 & 55 \\
Travessões anti-IA & 220 & 0 \\
Feedback de revisores & 10 & 0 \\
\bottomrule
\end{longtable}

\textbf{Inovações introduzidas:}
\begin{itemize}
    \item 5 revisores simulados com personas acadêmicas distintas;
    \item 4 orientadores PhD com perfis complementares;
    \item 6 engines de correção (gramatical, ABNT, anti-IA, referências, formatação, estilo);
    \item Conversão do AGENTS.md para chinês (-40\% tokens);
    \item Corretor CJK com 17 blocos Unicode e zero tolerância a vazamento.
\end{itemize}

% ======================================================================
% CAPÍTULO 4: SEGUNDA GERAÇÃO (EVO-7 A EVO-9)
% ======================================================================
\chapter{Segunda Geração: Sincronização Autônoma}

\section{Evo-7: Ecosystem Sync v3.5 + CJK Corrector (Score: 96/100)}

Marcou a transição de operação manual para sincronização autônoma:
\begin{itemize}
    \item \texttt{ecosystem-sync.ts} v3.5: scan periódico de 97 componentes;
    \item \texttt{manus-evolve.ts} v2.0: ciclo PLAN$\rightarrow$ACT$\rightarrow$REFLECT;
    \item \texttt{ptbr\_corrector.py}: 17 blocos Unicode CJK, 50+ regras ortográficas PT-BR;
    \item 210 arquivos com headers padronizados;
    \item 172 afinidades entre componentes;
    \item Matriz de afinidade computada automaticamente.
\end{itemize}

\section{Evo-8: Progressive Disclosure + Observability (Score: 98/100)}

\begin{itemize}
    \item 865 arquivos (+66\% em relação ao ciclo anterior);
    \item 76 skills, 10 diretórios, SKILL.md $<$ 2.5KB;
    \item Token budget por skill para eficiência de contexto;
    \item Agent Observability Monitor com Health Score;
    \item Cross-ecosystem affinity score: 95/100.
\end{itemize}

\textbf{Inovação:} Progressive Disclosure Design -- skills grandes são divididas entre SKILL.md (gatilhos) e \texttt{references/} (detalhes), otimizando o uso do contexto.

\section{Evo-9: Instalação em Massa de Repositórios (Score: 100/100)}

Expansão massiva via integração de 5 repositórios externos:
\begin{longtable}{@{}lcc@{}}
\caption{Repositórios Externos Integrados no Evo-9}\\
\toprule
\textbf{Repositório} & \textbf{Stars} & \textbf{Skills} \\
\midrule
addyosmani/agent-skills & 28.3k & 12 \\
nexu-io/open-design & 25.7k & 50+ \\
jimliu/baoyu-skills & 17.1k & 15+ \\
farmage/opencode-skills & -- & 10+ \\
cyijun/agent-smith & -- & 1 \\
\bottomrule
\end{longtable}

\textbf{Novos MCPs:} \texttt{mcp-pandoc}, \texttt{mcp-python-interpreter}, \texttt{mem0-mcp}.

\textbf{Resultado:} 209 skills combinadas, 20 MCPs ativos, score 100/100 pela primeira vez.

% ======================================================================
% CAPÍTULO 5: TERCEIRA GERAÇÃO (EVO-10 A EVO-13)
% ======================================================================
\chapter{Terceira Geração: Orquestração Multiagente}

\section{Evo-10: PyPI Scout + Ecosystem Hooks (Score: 95/100)}

\begin{itemize}
    \item \texttt{pypi\_scout.py} (350 linhas): catálogo curado de 22+ pacotes Python;
    \item Matriz de afinidade entre pacotes e 5 pipelines de integração;
    \item 5 Ecosystem Hooks v1.0 com verificação de compatibilidade;
    \item 7 bibliotecas instaladas: numpy, scipy, sympy, pandas, matplotlib, seaborn, plotly.
\end{itemize}

\section{Evo-11: DataOrchestrator (Score: 97/100)}

\begin{itemize}
    \item \texttt{data\_orchestrator.py} (592 linhas): consulta em linguagem natural $\to$ 8 domínios;
    \item 5 novos Ecosystem Hooks (GeoAnalyzer, FinanceAnalyzer, BioMedAnalyzer, etc.);
    \item 30+ bibliotecas instaladas, 20+ fontes Qualis A1;
    \item Integração com World Bank API, IBGE SIDRA, INPE, BCB, OpenAIRE.
\end{itemize}

\section{Evo-12: Sistema de Auditoria Caixa Branca (Score: 95/100)}

9 componentes de auditoria integrados:
\begin{enumerate}
    \item InteractionLogger -- log de todas as interações agente$\leftrightarrow$ferramenta;
    \item AcademicAuditTrail -- rastreabilidade de decisões acadêmicas;
    \item TokenEconomyMonitor -- custo por raciocínio, alerta de orçamento;
    \item AuditInstrumentor -- instrumentação automática de novos componentes;
    \item ResearcherScore -- score composto de qualidade;
    \item BudgetAlert -- notificação de estouro de orçamento;
    \item AuditDashboard -- painel visual de auditoria;
    \item AuditSearch -- busca semântica no histórico de auditoria;
    \item PipelineIntegration -- integração com o pipeline de evolução.
\end{enumerate}

\section{Evo-13: Reasoning Orchestrator v9.0 (Score: 96/100)}

\begin{itemize}
    \item 68 tipos de raciocínio (58 base + 10 Teoria dos Jogos);
    \item Bridge com o sistema de auditoria (AuditSystem);
    \item Integração com equilíbrio de Nash, Harsanyi e Shapley;
    \item Seleção adaptativa do framework lógico por tipo de problema.
\end{itemize}

% ======================================================================
% CAPÍTULO 6: ARQUITETURA CORA-DEBATE
% ======================================================================
\chapter{Arquitetura Cora-Debate}

\section{Visão Geral}

Cora (Cognitive ORchestrated Argumentation) é o núcleo de verificação formal do ecossistema. Combina 6 verificadores simbólicos (V1-V6) com um mecanismo de seleção adaptativa baseado em Q-Score UCB1.

\section{Verificadores Simbólicos}

\begin{longtable}{@{}lcc@{}}
\caption{Verificadores Cora-Debate (V1-V6)}\\
\toprule
\textbf{ID} & \textbf{Verificador} & \textbf{Motor} \\
\midrule
V1 & Análise Dimensional & Ontologia de unidades \\
V2 & Verificação Algébrica & SymPy simplify \\
V3 & Geração de Contraexemplos & Busca randomizada \\
V4 & Verificação Estatística & SciPy stats \\
V5 & Verificação Numérica & IEEE 754 float64 \\
V6 & Validação de EDOs/PDEs & SymPy dsolve \\
\bottomrule
\end{longtable}

\section{Q-Score UCB1}

O mecanismo de seleção de debatedores usa a fórmula:
\begin{equation}
    Q_i(N) = \bar{v}_i + \sqrt{\frac{2\ln N}{n_i}}
\end{equation}

onde $\bar{v}_i$ é o desempenho médio do debatedor $i$, $n_i$ é o número de vezes que foi selecionado e $N$ é o total de debates. O termo de exploração garante que debatedores menos testados sejam eventualmente avaliados, com regret bound $O(\log N)$.

\section{Resultados Experimentais}

\begin{longtable}{@{}lccc@{}}
\caption{Resultados do Cora-Debate}\\
\toprule
\textbf{Métrica} & \textbf{Inicial} & \textbf{Final} & \textbf{Ganho} \\
\midrule
Acurácia global & 65\% & 99\% & +34 pp ($p=3\times10^{-7}$) \\
Diversidade de raciocínio & 0.168 & 0.430 & +156\% \\
Problemas IMO resolvidos & -- & 55 & -- \\
PCI médio & -- & 98.3/100 & -- \\
Testes funcionais & 38 & 38 & 100\% \\
\bottomrule
\end{longtable}

\section{Especificação ANTISYM-001}

A SPEC-ANTISYM-001 detecta erros de antisimetria no produto exterior. O caso paradigmático é o erro clássico no artigo de Hénon-Heiles:

\begin{equation}
F^*(dx \wedge dy) = +b\;dx \wedge dy
\end{equation}

Quando o correto é:
\begin{equation}
F^*(dx \wedge dy) = \det(DF)\;dx \wedge dy = -b\;dx \wedge dy
\end{equation}

\textbf{Critérios de aceitação:}
\begin{enumerate}
    \item Detecta $+b\;dx\wedge dy$ em contexto de pullback;
    \item Aceita $-b\;dx\wedge dy$ como correto;
    \item Não dispara falso positivo para expressões su(2);
    \item Pontua $< 0.5$ para violações;
    \item Zero falso positivo para 3-formas.
\end{enumerate}

% ======================================================================
% CAPÍTULO 7: PIPELINE DE ENGENHARIA REVERSA
% ======================================================================
\chapter{Pipeline de Engenharia Reversa}

\section{Estrutura do Pipeline}

O pipeline Reversa executa engenharia reversa em 11 fases com 9 agentes especializados:

\begin{longtable}{@{}lll@{}}
\caption{Agentes do Pipeline Reversa}\\
\toprule
\textbf{Agente} & \textbf{Função} & \textbf{Estágio} \\
\midrule
reversa-scout & Mapeamento superficial & 1. Descoberta \\
reversa-archaeologist & Análise estática profunda & 2. Escavação \\
reversa-detective & Reconstrução lógica & 3. Dedução \\
reversa-architect & Arquitetura e ADRs & 4. Síntese \\
reversa-writer & Geração de SDDs & 5. Documentação \\
reversa-reviewer & Revisão cruzada & 6. Qualidade \\
reversa-visor & Visão sistêmica & 7. Integração \\
reversa-data-master & Modelagem de dados & 8. Dados \\
reversa-design-system & Sistema de design & 9. UI/UX \\
\bottomrule
\end{longtable}

\section{Métricas Obtidas}

\begin{itemize}
    \item Versão: 1.2.22, Confiança: 100/100 (anterior 87);
    \item 11/11 fases concluídas, 9/9 agentes executados;
    \item 67 artefatos criados, 12 gaps identificados e resolvidos;
    \item 12 ADRs registradas, 12 SDDs gerados;
    \item 12 checkpoints, 3 diagramas C4 (contexto, containers, componentes);
    \item 4 dias de duração do pipeline completo.
\end{itemize}

% ======================================================================
% CAPÍTULO 8: INTEGRAÇÃO MIROFISH/BETTAFISH
% ======================================================================
\chapter{Integração MiroFish/BettaFish: 18 Padrões Arquiteturais}

\section{Visão Geral}

O ecossistema incorporou 18 padrões arquiteturais (P1-P18) extraídos dos projetos MiroFish (666ghj/MiroFish), BettaFish (666ghj/BettaFish) e DeerFlow (bytedance/deer-flow).

\begin{longtable}{@{}lll@{}}
\caption{Padrões Arquiteturais MiroFish/BettaFish}\\
\toprule
\textbf{Padrão} & \textbf{Nome} & \textbf{Origem} \\
\midrule
P1 & Entity NER Reader & MiroFish \\
P2 & Hybrid Graph Retrieval & MiroFish \\
P3 & Graph Builder Pipeline & MiroFish \\
P4 & Ontology Generator & MiroFish \\
P5 & OASIS Profile Gen & MiroFish \\
P6 & Synthesis Agent (ReACT) & BettaFish \\
P7 & Swarm Review & OASIS \\
P8 & Code GraphRAG & MiroFish \\
P9 & Machine States & MiroFish \\
P10 & Graph Memory Updater & MiroFish \\
P11 & Process Lifecycle & MiroFish \\
P12 & Filesystem IPC & MiroFish \\
P13 & Config Generator & MiroFish \\
P14 & Agent Forum & BettaFish \\
P15 & Document IR & BettaFish \\
P16 & Agent Node Pipeline & BettaFish \\
P17 & MiddlewareChain + Reducers & DeerFlow \\
P18 & PhD Auditor & nexus-phd-strategist \\
\bottomrule
\end{longtable}

\section{P14: Agent Forum}

O Fórum de Agentes orquestra debates entre 125 agentes com 212 tipos de raciocínio em 27 categorias. O moderador OASIS controla turnos, evita repetições e grava o grafo de argumentos para análise posterior.

\begin{equation}
    \text{Pontuação}_{\text{Fórum}} = \frac{1}{N}\sum_{i=1}^N \left(w_1 \cdot \text{Qualidade}_i + w_2 \cdot \text{Diversidade}_i - w_3 \cdot \text{Repetição}_i\right)
\end{equation}

\section{P15: Document IR}

Sistema de Information Retrieval documental que acessa 50+ métricas reais (World Bank, WHO, FAO, UNESCO, SIPRI, ILO, ITU, IBGE, INPE) para fundamentar debates com dados verificáveis.

\section{P18: PhD Auditor}

Framework de auditoria acadêmica com 5 componentes:
\begin{itemize}
    \item \textbf{NashSolver:} Equilíbrio de Nash $N\times M$ para estratégias de debate;
    \item \textbf{StatisticalRigor:} Kappa de Cohen, correção de Bonferroni, Power Analysis ($1-\beta$);
    \item \textbf{QualisA1Auditor:} Score 0-100 com 10 critérios;
    \item \textbf{SensitivityAnalyzer:} Análise de sensibilidade paramétrica;
    \item \textbf{IMRADFormatter:} Formatação LaTeX/IMRAD com TSAC anti-AI.
\end{itemize}

% ======================================================================
% CAPÍTULO 9: TDD ACADÊMICO
% ======================================================================
\chapter{Sistema TDD Acadêmico: Validação de SPECs}

\section{Visão Geral}

O TDD Acadêmico converte cada recomendação da Meta-Avaliação DCA em uma Especificação Técnica (SPEC) com Critérios de Aceitação (CAs) implementados como testes automatizados. A arquitetura segue 4 design patterns:

\begin{itemize}
    \item \textbf{Template Method:} \texttt{SpecValidator.validar()} define o esqueleto;
    \item \textbf{Factory Method:} \texttt{TestRunner.criar\_padrao()} cria validadores;
    \item \textbf{Strategy:} cada SPEC é uma estratégia concreta;
    \item \textbf{Value Object:} \texttt{TestResult} e \texttt{SpecReport} são imutáveis.
\end{itemize}

\section{As 5 SPECs}

\begin{longtable}{@{}llll@{}}
\caption{Especificações Técnicas (SPECs)}\\
\toprule
\textbf{SPEC} & \textbf{Domínio} & \textbf{CAs} & \textbf{Score} \\
\midrule
PCI-001 & Calibração do Process Confidence Index & 5 & 100\% \\
CODE-001 & Verificador de Obrigatoriedade de Código & 5 & 100\% \\
ANTISYM-001 & Antisimetria do Produto Exterior ($\wedge$) & 5 & 100\% \\
NARR-001 & Conversor Narração $\to$ Demonstração & 5 & 100\% \\
CORA-001 & Expansão do Cora-Debate (V1-V6) & 5 & 100\% \\
\midrule
\textbf{Total} & -- & \textbf{25} & \textbf{100\%} \\
\bottomrule
\end{longtable}

\section{SPEC-PCI-001: Calibração do Process Confidence Index}

\textbf{Problema:} O PCI bruto superestima a confiança em artigos com simulações numéricas, especialmente em geometria e problemas numéricos.

\textbf{Fórmula de calibração:}
\begin{equation}
    \text{PCI}_{\text{cal}} = \min(10, \max(0, \text{PCI}_{\text{bruto}} \times 0.062 \times f_{\text{domínio}} + b_{\text{código}} + b_{\text{narração}}))
\end{equation}

\textbf{Critérios de aceitação:}
\begin{itemize}
    \item CA1: PCI 95 em geometria $\to$ $\approx$5.01;
    \item CA2: PCI nunca ultrapassa 10;
    \item CA3: PCI nunca é negativo;
    \item CA4: Domínio numérico penaliza mais que geral;
    \item CA5: Bônus por blocos de código presentes.
\end{itemize}

\section{SPEC-CODE-001: Verificador de Obrigatoriedade de Código}

\textbf{Problema:} Artigos que descrevem resultados numéricos sem disponibilizar código-fonte.

\textbf{Gatilhos de rejeição (9):} código, simulação, RK45, Runge-Kutta, Euler-Maruyama, Monte Carlo, ``o código confirma'', ``erro $<$ 10...'', ``resultado numérico''.

\textbf{Gatilhos de aviso (5):} dados, dataset, amostra, ``gráfico mostra'', ``figura mostra''.

\section{SPEC-NARR-001: Conversor Narração $\to$ Demonstração}

\textbf{Problema:} Uso excessivo de verbos narrativos que mascaram ausência de demonstração formal.

\textbf{Padrões detectados (N-01 a N-10):} obtém-se, verifica-se, deve-se, podemos, nota-se, afirma-se, é claro que, o código confirma, substituindo obtém-se, após simplificação.

\chapter{Resultados da Validação TDD}

Os 25 testes foram executados com 100\% de aprovação:

\begin{lstlisting}[caption=Saída do TDD Academic v2.0 (OOP)]
[PASS] PCI-001: 5/5 (100%) [##########]
[PASS] CODE-001: 5/5 (100%) [##########]
[PASS] ANTISYM-001: 5/5 (100%) [##########]
[PASS] NARR-001: 5/5 (100%) [##########]
[PASS] CORA-001: 5/5 (100%) [##########]
----------------------------------------------------------------
TOTAL: 25/25 (100.0%)
\end{lstlisting}

% ======================================================================
% CAPÍTULO 10: PhD AUDITOR E QUALIDADE QUALIS A1
% ======================================================================
\chapter{PhD Auditor e Qualidade Qualis A1}

\section{Critérios de Qualidade}

O PhD Auditor avalia a produção acadêmica em 10 critérios ponderados:

\begin{longtable}{@{}lcc@{}}
\caption{Critérios de Qualidade PhD Auditor}\\
\toprule
\textbf{Critério} & \textbf{Peso} & \textbf{Ideal} \\
\midrule
Densidade de evidências & 25\% & $\geq$ 2 evidências/parágrafo \\
Fontes verificadas (DOI) & 20\% & 100\% DOIs confirmados \\
TSAC compliance (87 palavras) & 20\% & 0 violações \\
Diversidade de fontes & 15\% & Fontes únicas $\geq$ parágrafos \\
Cobertura multi-domínio & 10\% & 8 domínios acessados \\
Peer review simulado & 10\% & 100\% parágrafos revisados \\
\bottomrule
\end{longtable}

\section{Escala de Qualidade}

\begin{itemize}
    \item \textbf{A ($\geq$90):} Qualis A1 -- padrão para periódicos de excelência;
    \item \textbf{B ($\geq$75):} Qualis A2-B1 -- qualidade superior;
    \item \textbf{C ($\geq$60):} Qualis B2-B3 -- qualidade mediana;
    \item \textbf{D ($\geq$40):} Qualis B4-C -- qualidade regular;
    \item \textbf{F ($<$40):} abaixo do padrão acadêmico mínimo.
\end{itemize}

\section{Resultados Alcançados}

\begin{longtable}{@{}lccc@{}}
\caption{Evolução do Auto Score Qualis}\\
\toprule
\textbf{Ciclo} & \textbf{Score} & \textbf{Classificação} & \textbf{Variação} \\
\midrule
Inicial & 74 & C & -- \\
Evo-6 (1ª correção) & 86.5 & B & +12.5 \\
Evo-6 (final) & 92.7 & A (Qualis A1) & +6.2 \\
Evo-7 (sync v3.5) & 95 & A (Qualis A1) & +2.3 \\
Evo-8 (progressive disclosure) & 98 & A (Qualis A1) & +3.0 \\
\bottomrule
\end{longtable}

% ======================================================================
% CAPÍTULO 11: MÉTRICAS AGREGADAS
% ======================================================================
\chapter{Métricas Agregadas do Ecossistema}

\section{Dimensões do Sistema}

\begin{longtable}{@{}lcr@{}}
\caption{Métricas Agregadas do Ecossistema OpenCode}\\
\toprule
\textbf{Dimensão} & \textbf{Métrica} & \textbf{Valor} \\
\midrule
\multirow{3}{*}{Agentes} & Total & 125 \\
 & Core & 56 \\
 & Criação (MASWOS) & 49 \\
 & SEEKER & 12 \\
 & Reversa & 7 \\
 & Corretor linguístico & 1 \\
\midrule
\multirow{3}{*}{MCPs} & Configurados & 40 \\
 & Ativos & 17 \\
 & Locais/Remotos & 38/2 \\
\midrule
\multirow{2}{*}{Skills} & Total & 106 \\
 & Categorias & 13 \\
\midrule
\multirow{2}{*}{Plugins} & npm & 10 \\
 & TypeScript locais & 2 \\
 & Bridge & 3 \\
\midrule
\multirow{3}{*}{Raciocínio} & Tipos totais & 212 \\
 & Categorias & 27 \\
 & Teoria dos Jogos & 10 \\
\midrule
\multirow{2}{*}{Evolução} & Ciclos & 13 \\
 & Score trajetória & 85 $\to$ 100 \\
\midrule
\multirow{2}{*}{Código} & Linhas Python & $\sim$120.000 \\
 & Arquivos & $>$ 865 \\
\midrule
\multirow{2}{*}{Qualidade} & Health Score & 100/100 \\
 & Testes TDD & 25/25 \\
\midrule
\multirow{2}{*}{Acadêmico} & Fontes Qualis A1 & 20+ \\
 & Editais curados & 52 \\
\bottomrule
\end{longtable}

\section{Evolução do Health Score}

\begin{longtable}{@{}lcc@{}}
\caption{Evolução do Health Score ao Longo dos Ciclos}\\
\toprule
\textbf{Data} & \textbf{Score} & \textbf{Otimização} \\
\midrule
05/05 22:45 & 85 & Inicial \\
06/05 22:20 & 90 & sync-v3.5-initial \\
06/05 22:27 & 96 & sync-v3.5-final \\
06/05 23:35 & 100 & sync-v4.0-complete \\
07/05 00:00 & 100 & artigo-qualis-a1-95 \\
08/05 18:00 & 100 & progressive-disclosure-round8 \\
09/05 12:00 & 100 & round9-install-evolve \\
12/05 19:30 & 100 & round10-mcpick-sdk-update \\
14/05 17:00 & 100 & round11-academic-ml-export-abnt \\
14/05 20:34 & 100 & round12-artigo-v3-complexidade \\
\bottomrule
\end{longtable}

% ======================================================================
% CAPÍTULO 12: CONCLUSÃO
% ======================================================================
\chapter{Conclusão}

\section{Contribuições}

Este trabalho documentou a evolução completa do Ecossistema OpenCode ao longo de 13 ciclos evolutivos, demonstrando que:

\begin{enumerate}
    \item \textbf{É possível evolução autônoma:} o ciclo PLAN$\rightarrow$ACT$\rightarrow$REFLECT$\rightarrow$EXTRACT$\rightarrow$EVOLVE permitiu que o sistema evoluísse de score 85 para 100 sem intervenção humana;
    \item \textbf{Verificação formal é viável em IA:} o Cora-Debate com 6 verificadores (V1-V6) alcançou 99\% de acurácia e 100\% nos testes funcionais;
    \item \textbf{Integração multi-repositório funciona:} 18 padrões arquiteturais (P1-P18) de 3 projetos upstream foram integrados com sucesso;
    \item \textbf{TDD acadêmico garante qualidade:} 5 SPECs com 25 critérios de aceitação validados a 100\%;
    \item \textbf{Padrão Qualis A1 é alcançável:} o sistema atingiu score 98/100 nos critérios PhD Auditor.
\end{enumerate}

\section{Trabalhos Futuros}

\begin{enumerate}
    \item Expansão para 300+ tipos de raciocínio em 35+ categorias;
    \item Integração com modelos de linguagem multimodal (visão + texto);
    \item Publicação do ecossistema como plataforma open-source;
    \item Validação em domínios clínicos (aplicação em arteterapia);
    \item Expansão do Cora-Debate para verificação de código-fonte (V7).
\end{enumerate}

% ======================================================================
% REFERÊNCIAS
% ======================================================================
\newpage
\bibliographystyle{abntex2-alf}
\bibliography{dissertacao_evolucao_opencode}

% ======================================================================
% APÊNDICES
% ======================================================================
\begin{apendicesenv}

\chapter{Relatório TDD Completo}

\section{Resultados por SPEC}

\begin{longtable}{@{}lll@{}}
\toprule
\textbf{SPEC} & \textbf{Resultado} & \textbf{Taxa} \\
\midrule
PCI-001 & 5/5 PASS & 100\% \\
CODE-001 & 5/5 PASS & 100\% \\
ANTISYM-001 & 5/5 PASS & 100\% \\
NARR-001 & 5/5 PASS & 100\% \\
CORA-001 & 5/5 PASS & 100\% \\
\midrule
\textbf{Total} & \textbf{25/25 PASS} & \textbf{100\%} \\
\bottomrule
\end{longtable}

\section{Detalhamento PCI-001}

\begin{itemize}
    \item [+] CA1: PCI bruto 95, geometria, 0 blocos $\to$ $\approx$5.01
    \item [+] CA2: PCI calibrado nunca ultrapassa 10.0
    \item [+] CA3: PCI calibrado nunca é negativo
    \item [+] CA4: Domínio ``numérico'' sempre penaliza mais que ``geral''
    \item [+] CA5: Bônus por código só se bloco REAL existe
\end{itemize}

\section{Detalhamento CODE-001}

\begin{itemize}
    \item [+] CA1: ``Código Python com RK45'' sem bloco $\to$ rejeitado
    \item [+] CA2: Solução puramente analítica sem gatilhos $\to$ aprovado
    \item [+] CA3: Código em arquivo externo $\to$ aprovado
    \item [+] CA4: Falso positivo zero para texto matemático puro
    \item [+] CA5: Mensagem de rejeição aponta o gatilho específico
\end{itemize}

\section{Detalhamento ANTISYM-001}

\begin{itemize}
    \item [+] CA1: Detecta $+b\;dx\wedge dy$ onde deveria ser $-b\;dx\wedge dy$
    \item [+] CA2: Não dispara falso positivo para pullback correto
    \item [+] CA3: Não dispara falso positivo para su(2)
    \item [+] CA4: Pontuação $<$ 0.5 dispara alerta
    \item [+] CA5: Zero falso positivo em $n$-formas com $n>1$
\end{itemize}

\section{Detalhamento NARR-001}

\begin{itemize}
    \item [+] CA1: 100\% dos padrões N-01 a N-10 são detectados
    \item [+] CA2: ``o código confirma'' sem bloco $\to$ pendente
    \item [+] CA3: Expressões ``obtém-se'' são detectáveis
    \item [+] CA4: ``é claro que'' sempre dispara verificação
    \item [+] CA5: Taxa de narração calculada corretamente
\end{itemize}

\section{Detalhamento CORA-001}

\begin{itemize}
    \item [+] CA1: V4 detecta $\{J_1,J_2\}=+J_0$ em contexto su(1,1)
    \item [+] CA2: V4 aceita $[J_1,J_2]=+J_3$ em contexto su(2)
    \item [+] CA3: V4 rejeita $[H,E]=-2E$ onde deveria ser $+2E$ (sl(2,R))
    \item [+] CA4: V5 sugere álgebra alternativa (su(2) para $[J_1,J_2]=+J_3$)
    \item [+] CA5: V6 verifica conectivos lógicos (logo, portanto)
\end{itemize}

% ======================================================================
% APÊNDICE B: DIAGRAMA DE CLASSES TDD
% ======================================================================
\chapter{Arquitetura de Classes do TDD Acadêmico}

\begin{verbatim}
SpecValidator (ABC)
|   +validar() -> SpecReport  (Template Method)
|   #_executar_testes() -> List[TestResult]  (abstract)
|
+-- PciValidator      (SPEC-PCI-001)
+-- CodeValidator     (SPEC-CODE-001)
+-- AntisymValidator  (SPEC-ANTISYM-001)
+-- NarrValidator     (SPEC-NARR-001)
+-- CoraValidator     (SPEC-CORA-001)

TestRunner (Composicao)
|   +criar_padrao() -> TestRunner  (Factory Method)
|   +executar_todos() -> List[SpecReport]

TestReport
|   +exibir() -> stdout
|   +to_json(path) -> JSON
|   +to_markdown() -> MD

TestResult (Value Object)
|   -spec, -ca, -descricao, -status, -erro
|   +aprovou() -> bool

SpecReport (Value Object)
|   -spec_id, -resultados, -total, -passes, -taxa
\end{verbatim}

% ======================================================================
% APÊNDICE C: GLOSSÁRIO
% ======================================================================
\chapter{Glossário de Termos Técnicos}

\begin{longtable}{@{}lp{10cm}@{}}
\toprule
\textbf{Termo} & \textbf{Definição} \\
\midrule
ADR & Architecture Decision Record -- registro formal de decisão arquitetural \\
ANP & Agent Node Pipeline -- pipeline de nós tipados e composáveis \\
CA & Critério de Aceitação -- condição que um teste deve satisfazer \\
Cora & Cognitive ORchestrated Argumentation -- arquitetura de verificação formal \\
DCA & Dinâmica Caótica Aplicada -- domínio de aplicação principal \\
IPC & Inter-Process Communication -- comunicação entre processos \\
MCP & Model Context Protocol -- protocolo de comunicação entre agente e ferramenta \\
NER & Named Entity Recognition -- extração de entidades nomeadas \\
OASIS & Open Agent Simulation and Integration System -- simulador multiagente \\
PCI & Process Confidence Index -- índice de confiança do processo \\
Q-Score & Quality Score -- pontuação de qualidade com UCB1 \\
SDD & Software Design Document -- documento de design de software \\
SPEC & Especificação Técnica -- documento com critérios de aceitação \\
TDD & Test-Driven Development -- desenvolvimento orientado a testes \\
TSAC & Textual Source Analysis Credibility -- sistema de credibilidade de fontes \\
UCB1 & Upper Confidence Bound -- algoritmo de seleção adaptativa \\
V1-V6 & Verificadores simbólicos do Cora-Debate \\
\bottomrule
\end{longtable}

\end{apendicesenv}

\end{document}
